\changecolours{ScoutGreen}
\section[Exercícios]{Exercícios e Desafios Práticos Resolvidos}

% ------------------------------------------------------------------------------
\twocolframe[title={Desafio 1: Biblioteca e Empréstimos por Categoria},leftcol=7.0cm,rightcol=7.0cm]{%
  \alert{Enunciado do Problema}\par
  \itemseps[topsepi=0.5mm,itemsepi=1mm,labelsep=2mm,leftmargini=3mm]
  \begin{itemize}
    \item Recupere o \texttt{titulo} da obra, o nome do \texttt{leitor} e a \texttt{data\_emp} dos empréstimos realizados para livros da categoria \textbf{'Computacao'}.
    \item Esquemas das tabelas:
      \begin{itemize}
        \item \texttt{LIVRO(isbn, titulo, categoria, ano)}
        \item \texttt{EMPRESTIMO(id\_emp, isbn, leitor, data\_emp)}
      \end{itemize}
  \end{itemize}
}{%
  \alert{Resolução Formal}\par
  \itemseps[topsepi=0.5mm,itemsepi=1mm,labelsep=2mm,leftmargini=3mm]
  \begin{itemize}
    \item \textbf{Estratégia Otimizada (Push-Down Selection):}
      {\footnotesize
      \begin{align*}
        \text{COMP} &\leftarrow \sigma_{\text{categoria} = '\text{Computacao}'}(\text{LIVRO}) \\
        \text{JUNCAO} &\leftarrow \text{COMP} \bowtie \text{EMPRESTIMO} \\
        \text{RES} &\leftarrow \pi_{\text{titulo, leitor, data\_emp}}(\text{JUNCAO})
      \end{align*}}
    \item \textbf{Tuplas Obtidas:}
      \begin{itemize}
        \item $\langle \text{'Algoritmos'}, \text{'Alice'}, \text{'10/11'} \rangle$
        \item $\langle \text{'Algoritmos'}, \text{'Bob'}, \text{'15/11'} \rangle$
      \end{itemize}
  \end{itemize}
}

% ------------------------------------------------------------------------------
\twocolframe[title={Desafio 2: Livros Nunca Emprestados via Anti-Junção},leftcol=7.0cm,rightcol=7.0cm]{%
  \alert{Enunciado do Problema}\par
  \itemseps[topsepi=0.5mm,itemsepi=1mm,labelsep=2mm,leftmargini=3mm]
  \begin{itemize}
    \item Identifique o \texttt{isbn} e o \texttt{titulo} de todos os livros do acervo que \textbf{nunca foram emprestados} a nenhum leitor.
    \item Requisito: Não utilize produtos cartesianos ou operações arriscadas de memória.
  \end{itemize}
}{%
  \alert{Resolução com Anti-Junção}\par
  \itemseps[topsepi=0.5mm,itemsepi=1mm,labelsep=2mm,leftmargini=3mm]
  \begin{itemize}
    \item \textbf{Aplicação Direta de $\triangleright$:}
      {\footnotesize
      \begin{align*}
        \text{SEM\_EMP} &\leftarrow \text{LIVRO} \triangleright \text{EMPRESTIMO} \\
        \text{RES} &\leftarrow \pi_{\text{isbn, titulo}}(\text{SEM\_EMP})
      \end{align*}}
    \item \textbf{Em Operadores Fundamentais:}
      \[ \pi_{\text{isbn, titulo}}\left(\text{LIVRO} - \pi_{\text{atr(LIVRO)}}(\text{LIVRO} \bowtie \text{EMPRESTIMO})\right) \]
    \item \textbf{Tuplas Obtidas:}
      \begin{itemize}
        \item $\langle \text{'L03'}, \text{'Quimica Geral'} \rangle$
      \end{itemize}
  \end{itemize}
}

% ------------------------------------------------------------------------------
\twocolframe[title={Desafio 3: Relatório de Vendas com Left Outer Join},leftcol=7.0cm,rightcol=7.0cm]{%
  \alert{Enunciado do Problema}\par
  \itemseps[topsepi=0.5mm,itemsepi=1mm,labelsep=2mm,leftmargini=3mm]
  \begin{itemize}
    \item Gere uma listagem com \textbf{todos} os clientes cadastrados exibindo \texttt{nome}, \texttt{cidade} e o \texttt{id\_ped} correspondente, garantindo que os clientes que ainda não compraram não sejam descartados.
  \end{itemize}
}{%
  \alert{Resolução Formal}\par
  \itemseps[topsepi=0.5mm,itemsepi=1mm,labelsep=2mm,leftmargini=3mm]
  \begin{itemize}
    \item \textbf{Junção Externa Esquerda:}
      \[ \text{RES} \leftarrow \pi_{\text{nome, cidade, id\_ped}}(\text{CLIENTE} \mathrel{\sqsubset\!\!\bowtie} \text{PEDIDO}) \]
    \item \textbf{Instância Resultante:}
      {\tiny
      \begin{tabular}{llr}
        \toprule
        \textbf{nome} & \textbf{cidade} & \textbf{id\_ped} \\
        \midrule
        Ana Lima & Tauá & 501 \\
        Ana Lima & Tauá & 503 \\
        Bruno Silva & Fortaleza & \textit{null} \\
        Carla Dias & Tauá & 502 \\
        Diego Melo & Sobral & 504 \\
        Eva Rocha & Tauá & \textit{null} \\
        \bottomrule
      \end{tabular}}
  \end{itemize}
}

% ------------------------------------------------------------------------------
\twocolframe[title={Desafio 4: Conclusão de Núcleo Básico com Divisão},leftcol=7.0cm,rightcol=7.0cm]{%
  \alert{Enunciado do Problema}\par
  \itemseps[topsepi=0.5mm,itemsepi=1mm,labelsep=2mm,leftmargini=3mm]
  \begin{itemize}
    \item Recupere o \texttt{id\_alu} e o \texttt{nome} dos estudantes que cursaram e foram aprovados em \textbf{todas} as matérias contidas na tabela \texttt{DISC\_OBRIGATORIAS}.
    \item Esquemas: 
      \begin{itemize}
        \item \texttt{HISTORICO(id\_alu, cod\_disc)}
        \item \texttt{DISC\_OBRIGATORIAS(cod\_disc)}
        \item \texttt{ALUNO(id\_alu, nome, cidade)}
      \end{itemize}
  \end{itemize}
}{%
  \alert{Resolução Formal com Divisão}\par
  \itemseps[topsepi=0.5mm,itemsepi=1mm,labelsep=2mm,leftmargini=3mm]
  \begin{itemize}
    \item \textbf{Etapa 1 (Divisão):}
      \[ \text{APROVADOS} \leftarrow \text{HISTORICO} \div \text{DISC\_OBRIGATORIAS} \]
    \item \textbf{Etapa 2 (Recuperação do Nome por Junção):}
      \[ \text{RES} \leftarrow \pi_{\text{id\_alu, nome}}(\text{APROVADOS} \bowtie \text{ALUNO}) \]
    \item \textbf{Alunos Qualificados:}
      \begin{itemize}
        \item $\langle 101, \text{'Ana Lima'} \rangle$
        \item $\langle 104, \text{'Diana Moreira'} \rangle$
      \end{itemize}
  \end{itemize}
}
